% ---------------------------------------------------------------------------
% Subsection: verification of the coupled flow-transport scheme by MMS.
% Placeholders to resolve against the main document:
%   \ref{sec:flow}, \ref{sec:transport}, \ref{sec:coupling},
%   \ref{eq:cancelled}, \ref{tab:vgm}, \cite{roache2002,salari2000}
% New references (BibTeX in mms_coupled_refs.bib):
%   \cite{radu2014,suciu2021,shunn2012}
% Figures (figures/, made by mms_plot.py): mms_fields.pdf, mms_profiles.pdf,
%   mms_convergence.pdf
% ---------------------------------------------------------------------------
\subsection{Verification of the coupled scheme with manufactured solutions}
\label{sec:mms}

The coupling of Section~\ref{sec:coupling} is conservative and bounded by
construction, but neither property says anything about accuracy. We therefore
measure the convergence of the coupled discretization with the method of
manufactured solutions (MMS) \cite{roache2002,salari2000}. Because the
transport model advects with the fluxes produced by the flow model, and the
flow model in turn uses the density and viscosity of the transported
concentration, an error seen in the coupled solution can originate in either
model or in the exchange between them. We separate these contributions with a
hierarchy of four tests of increasing coupling strength, followed by a sweep
of the coupling strength itself.

\begin{enumerate}
  \item[\textbf{T0}] \emph{Flow only.} The flow model of
    Section~\ref{sec:flow} alone, with $\rho\equiv\rho_0$.
  \item[\textbf{T1}] \emph{Constant concentration.} The coupled model with
    $c\equiv c_0$ in the initial and boundary data and $\epsilon\neq0$. The
    exact concentration is $c=c_0$ and the flow problem reduces to T0 with
    gravity scaled by $\rho(c_0)/\rho_0$ and $\mathbf{K}_s$ by
    $\mu_0/\mu(c_0)$. We require (i) $c_h\equiv c_0$ to round-off and (ii)
    $\psi_h$ equal to the correspondingly scaled T0 solution to within the
    nonlinear-solver tolerance. The cancelled form \eqref{eq:cancelled}
    preserves constants for any set of fluxes, so T1 does not test transport
    accuracy; it tests the exchange of $\rho$, $\mu$ and the carried fluid
    mass.
  \item[\textbf{T2}] \emph{One-way coupling.} $\epsilon=0$ and constant
    viscosity, so the flow is independent of $c$ and the transport model is
    driven by the discrete fluxes of a flow solution that does not see it.
    The flow part of T2 is therefore T0.
  \item[\textbf{T3}] \emph{Two-way coupling.} The full model, with $\rho(c)$
    and $\mu(c)$ fed back to the flow model.
\end{enumerate}

T2 and T3 use the same manufactured fields, so comparing them shows how much
the feedback of $c$ on the flow, which the scheme lags by one time step
(Section~\ref{sec:coupling}), changes the errors.

\paragraph{Manufactured solution.}
Let $\Omega=(0,L)^2$ with $z$ pointing upward, $L=1$~m and $t\in(0,T]$. The
pressure head follows the manufactured solution of Radu and
Wang~\cite{radu2014}, which Suciu et al.~\cite{suciu2021} also use to verify a
fully coupled Richards--transport solver,
\begin{equation}
  \psi_e(x,z,t) = \psi_0\Big(1 + A_\psi\,\frac{t}{T}\,b(x,z)\Big),
  \qquad b(x,z) = 16\,\frac{x}{L}\Big(1-\frac{x}{L}\Big)\frac{z}{L}\Big(1-\frac{z}{L}\Big).
  \label{eq:mms_psi}
\end{equation}
With $\psi_0=-1$, $A_\psi=1/16$ and $L=T=1$, \eqref{eq:mms_psi} is the field of
\cite{radu2014,suciu2021}. We take $\psi_0=-1/\alpha$ and $A_\psi=1$, so that
$\psi_e\in[-2/\alpha,-1/\alpha]$ and the effective saturation spans a range wide
enough to exercise the nonlinearity of $k_r$. The concentration is adapted from
Problem~2 of Shunn et al.~\cite{shunn2012}, a front that advects at a constant
speed, widens in time and carries a transverse corrugation:
\begin{equation}
  c_e(x,z,t) = \bar c + A_c \tanh\!\left(\frac{e^{-\omega t}\,\xi(x,z,t)}{\delta}\right),
  \qquad
  \xi = z - z_f(t) + a_f\cos(2\pi x/L),
  \qquad z_f(t) = z_0 - U t.
  \label{eq:mms_c}
\end{equation}
In \cite{shunn2012} the transported scalar is the rational function of
$\tanh$ that makes the density a $\tanh$ profile under their mixing law. Our
density is linear in $c$, so $c_e$ itself is the $\tanh$ profile and $\rho$
follows it; we also align the front normal with gravity and drop the
transverse drift. The parameters are $\bar c=0.5$, $A_c=0.3$, $z_0=0.75L$,
$U=0.5L/T$, $\delta=0.15L$, $a_f=0.05L$ and $\omega=\ln 2/T$, so the front
width doubles over the run, together with the van Genuchten--Mualem parameters
of Table~\ref{tab:vgm}. The final time is one transit time,
$T=\theta_{\min}L/\max|\mathbf{v}|$, evaluated once for T3 with
$\epsilon=0.235$ ($T=873$~s), so that every test and every point of the
coupling sweep below shares the same $T$ and the same time steps.
The fields are chosen so that the test measures the discretization in its
smooth regime and nothing else:
\begin{itemize}
  \item $\psi_e\in[-2/\alpha,-1/\alpha]$, so the domain stays unsaturated
    ($S_e\in[0.35,0.67]$ for the parameters of Table~\ref{tab:vgm}); no
    seepage face activates and the retention and conductivity functions are
    smooth. Since $b=0$ on $\partial\Omega$, the head is the constant
    $\psi_0$ on the whole boundary and at $t=0$.
  \item $\partial_z\psi_e\ge-4A_\psi/(\alpha L)$ and $\rho/\rho_0\ge1$, so
    $v_z=-k_r K_{s,zz}\mu_0/\mu\,(\partial_z\psi_e+\rho/\rho_0)<0$ everywhere
    whenever $A_\psi<\alpha L/4$. The velocity never vanishes, and the
    dispersion tensor, which is not differentiable at $\mathbf{v}=0$, is
    smooth.
  \item $c_e$ is strictly monotone in $z$ and stays inside $(0.2,0.8)$, so it
    has no interior extrema and does not approach the bounds $[0,1]$. The FCT
    limiters then act only through the local bounds of a monotone field, and
    the observed rates are not reduced by clipping at smooth extrema.
\end{itemize}
The coupling differs from that of \cite{suciu2021}, where the concentration
enters the water content, $\theta(\psi,c)$: here it enters through
$\rho(c)$ in the storage and the buoyancy term and through $\mu(c)$ in the
conductivity, which are the coupling paths of Section~\ref{sec:coupling}. As a
cross-check, T3 is repeated with the concentration of \cite{suciu2021},
rescaled to $c=0.2+0.6\,(t/T)\,b(x,z)$, whose interior maximum the FCT
limiters must clip.

\paragraph{Source terms.}
Substituting $(\psi_e,c_e)$ into the flow equation and the conservative
transport equation, with $\rho=\rho(c_e,\psi_e)$ and $\mu=\mu(c_e)$, gives
\begin{align}
  f_w &= \partial_t(\rho\theta) + \nabla\cdot(\rho\mathbf{v}),
  \label{eq:mms_fw}\\
  f_s &= \partial_t(\rho\theta c) + \nabla\cdot(\rho\mathbf{v}c - \rho\theta\mathbf{D}\nabla c),
  \label{eq:mms_fs}
\end{align}
evaluated symbolically. The manufactured solutions of \cite{shunn2012} satisfy
the continuity equation without a source, because continuity is central to
their solution algorithm. The same holds for the coupling here, but the Darcy
velocity is fixed by $\psi$, and no closed-form pair $(\psi_e,c_e)$ with
$f_w=0$ is available for the van Genuchten--Mualem relations; we therefore
keep $f_w\neq0$ and test the exchange of the fluid balance separately in T1.
The transport model solves the cancelled form \eqref{eq:cancelled}, obtained
by subtracting $c$ times the flow equation from the transport equation. Its
source is therefore
\begin{equation}
  g = f_s - c\,f_w
    = \rho\theta\,\partial_t c + \rho\mathbf{v}\cdot\nabla c
      - \nabla\cdot(\rho\theta\mathbf{D}\nabla c).
  \label{eq:mms_g}
\end{equation}
Both sources are added as lumped nodal terms evaluated at $t^{n+1}$:
$-|\Omega_i|\,f_w(\mathbf{x}_i,t^{n+1})$ in the flow residual and
$-|\Omega_i|\,g(\mathbf{x}_i,t^{n+1})$ in the transport residual, where
$|\Omega_i|=\int_\Omega w_i\,d\mathbf{x}$. In the discrete transport equation
the fluid source reaches node $i$ through the nodal exchange $B_i$, which at
an interior node carries the local concentration $c_i$ and so drops out of
\eqref{eq:cancelled}. The solute added at node $i$ is thus
$|\Omega_i|(g+c_i f_w)$, consistent with \eqref{eq:mms_fs}. Passing $f_s$
instead of $g$ would leave an $O(1)$ error of size $c f_w$.

Each test is run first with a constant diffusion coefficient,
$\mathbf{D}=D_m\mathbf{I}$, as in \cite{suciu2021}, and then with mechanical
dispersion added, $\alpha_L=10^{-3}$~m and $\alpha_T=0.1\,\alpha_L$. We take
$D_m=10^{-6}$~m$^2$/s, the largest value for which the mesh P\'eclet number
$|\mathbf{v}|h/(\theta D)$ at the front stays above one on every mesh (from
23 at $N=10$ to 1.4 at $N=160$), so the transport is advection dominated
throughout the ladder. Dirichlet data $\psi=\psi_e$ (imposed weakly,
Section~\ref{sec:flow}) and $c=c_e$ are prescribed on all of $\partial\Omega$,
and the initial data are the nodal interpolants of $\psi_e(\cdot,0)$ and
$c_e(\cdot,0)$.

\paragraph{Refinement and error measures.}
We use structured triangulations with $N=10\cdot2^k$ intervals per side,
$k=0,\dots,4$, $h=L/N$, and refine space and time together, $\Delta t\propto h$
at a fixed Courant number
$\mathrm{Co}=\Delta t\,\max|\mathbf{v}|/(\theta_{\min}h)=0.5$. At $t=T/4$,
$T/2$ and $T$ we report the lumped-mass weighted discrete $L^2$ error
\begin{equation}
  \|e\|_{L^2}(c) = \Big(\sum_i |\Omega_i|\,\big|c_i - c_e(\mathbf{x}_i,t)\big|^2\Big)^{1/2},
\end{equation}
the same norm for $\psi$, and the observed order
$\log_2\big(\|e\|^{(k-1)}_{L^2}/\|e\|^{(k)}_{L^2}\big)$. Every run also reports the fluid continuity
defect of the transport model, i.e.\ the change in the carried fluid mass minus
the time-integrated exchange with the flow model, normalized by the initial
mass. By the construction of Section~\ref{sec:coupling} it should sit at the
level of the nonlinear-solver tolerance on every mesh, independently of the
discretization error.

Each test is run with both models in their low-order form (FCT off) and with
FCT on in both models, the configuration used in the remainder of this paper.
Both models use backward Euler, and the coupling lags $\rho(c)$ and $\mu(c)$ by
one step, i.e.\ the coupled equations are solved with a single outer iteration
per step, which limits the temporal order to one \cite{shunn2012}. Since
$\Delta t$ is refined with $h$, the observed orders combine the spatial and
temporal errors and tend to one as the mesh is refined; a rate above one
means that the spatial error, which converges faster, still dominates
\cite{shunn2012}.

\paragraph{Results.}
T1 is reproduced exactly: $c_h\equiv c_0$ at every node and step, and the
coupled $\psi_h$ differs from the correspondingly scaled flow-only solution by
at most $3\cdot10^{-16}$~m on $N=10$ and $N=20$. Figure~\ref{fig:mms_fields}
compares the manufactured and the computed fields of T3 at $t=T$ on $N=80$.
The concentration error, at most $4\cdot10^{-3}$, sits at the front: negative
above its centre and positive below it, the signature of numerical diffusion.
The head error, at most 0.43~cm against an 8.6~cm range of $\psi_e$, is largest
in the upper half of the domain, above the dry centre.
Figure~\ref{fig:mms_profiles} shows the same comparison along the vertical
line $x=L/2$ on $N=20$ and $N=80$.

Tables~\ref{tab:mms_t3_psi} and~\ref{tab:mms_t3_c} give the $L^2$ errors of
T3 at $t=T/4$, $T/2$ and $T$, and Figure~\ref{fig:mms_convergence} plots them
at $t=T$. The head is where the high-order flow flux pays off. With FCT, the
order of $\|e\|_{L^2}(\psi)$ is 1.1 to 1.2 at $t=T/4$, 1.2 to 1.3 at $t=T/2$
and reaches 1.51 at $t=T$, and the error on $N=160$ is 1.5, 2.5 and 3.3 times
smaller than with the low-order scheme at the three times. The low-order
scheme stays at or below first order, with orders of 0.77 to 0.99. An order
above one means that the spatial error, which converges faster, still dominates the
combined error, as discussed above.

The concentration converges at first order with FCT at every time, with
orders of 1.0 to 1.2. FCT reduces the concentration error on the coarse and
intermediate meshes, by 39 to 58\% at $t=T/4$ on $N\le40$, but the low-order
error converges faster on the finest meshes, with orders of 1.2 to 1.4, and on
$N=160$ the two schemes agree to within 6\% at $t=T$. On these meshes the
backward Euler error, which both schemes share, dominates the concentration
error.

T2 shows the same behaviour, with orders of $\|e\|_{L^2}(\psi)$ at $t=T$ of
1.57 (FCT) and 0.94 (low order) on $N=160$; since the flow of T2 does not see
$c$, these are the rates of T0. Adding mechanical dispersion changes the FCT
errors by at most 3.5\% and leaves the head unchanged; the low-order
concentration errors drop by up to 16\%. For the concentration of
\cite{suciu2021}, whose maximum is clipped by the limiters,
$\|e\|_{L^2}(c)$ still converges at 1.67 on $N=160$. The fluid continuity
defect stays below $1.4\cdot10^{-18}$ in every run.

\paragraph{Coupling strength.}
The only coupling error of the scheme is the one-step lag of $\rho(c)$ and
$\mu(c)$. We repeat T3 with FCT on $N=40$ for
$\epsilon\in\{0,0.1,0.235,0.4\}$ with constant viscosity, and for two
viscosity contrasts at $\epsilon=0.235$, keeping the manufactured fields and
recomputing the sources for each value (Table~\ref{tab:mms_sweep}). Over this
range $\|e\|_{L^2}(c)$ at $t=T$ changes by less than 10\%, and halving
$\Delta t$ reduces it by a factor of 1.70 to 1.96, so on this mesh the
concentration error is first order in time for every coupling strength. The
difference of the errors at $\epsilon$ and at $\epsilon=0$ is not a measure of
the lag error, since the sources, and with them the errors, change with
$\epsilon$; it is small compared with the errors themselves and changes sign,
so the feedback adds no error that the norms resolve.

\begin{figure*}[t]
  \centering
  \includegraphics[width=\textwidth]{mms_fields}
  \caption{Two-way coupling (T3), FCT, $N=80$, $t=T$: manufactured (left)
    and computed (middle) concentration and pressure head, and their
    difference (right).}
  \label{fig:mms_fields}
\end{figure*}

\begin{figure}[t]
  \centering
  \includegraphics[width=\columnwidth]{mms_profiles}
  \caption{Profiles of $c$ and $\psi$ along $x=L/2$ at $t=T$ for T3 with FCT:
    manufactured solution and computed solutions on $N=20$ and $N=80$.}
  \label{fig:mms_profiles}
\end{figure}

\begin{figure}[t]
  \centering
  \includegraphics[width=\columnwidth]{mms_convergence}
  \caption{Two-way coupling (T3): $\|e\|_{L^2}$ of $c$ and $\psi$ at $t=T$ against
    the mesh size, with $\Delta t\propto h$ at $\mathrm{Co}=0.5$. Dashed and
    dotted lines have slopes one and two.}
  \label{fig:mms_convergence}
\end{figure}

\begin{table}[t]
\centering
\begin{tabular}{c|cc|cc}
\hline
\multirow{2}{*}{$h_e$ [m]}
& \multicolumn{2}{c|}{\textbf{Low Order}}
& \multicolumn{2}{c}{\textbf{FCT}} \\
\cline{2-5}
& $\|e\|_{L^2}$ & Order
& $\|e\|_{L^2}$ & Order \\
\hline
\hline

\multicolumn{5}{c}{\textbf{$t = T/4 = 218.2$ s}} \\
\hline
$0.1$        & $1.641\times10^{-3}$ & --
           & $2.000\times10^{-3}$ & -- \\
$0.05$       & $8.302\times10^{-4}$ & $0.983$
           & $8.835\times10^{-4}$ & $1.179$ \\
$0.025$      & $4.174\times10^{-4}$ & $0.992$
           & $3.829\times10^{-4}$ & $1.206$ \\
$0.0125$     & $2.104\times10^{-4}$ & $0.988$
           & $1.674\times10^{-4}$ & $1.194$ \\
$0.00625$    & $1.136\times10^{-4}$ & $0.889$
           & $7.619\times10^{-5}$ & $1.135$ \\
\hline\hline

\multicolumn{5}{c}{\textbf{$t = T/2 = 436.4$ s}} \\
\hline
$0.1$        & $5.077\times10^{-3}$ & --
           & $4.897\times10^{-3}$ & -- \\
$0.05$       & $2.687\times10^{-3}$ & $0.918$
           & $2.181\times10^{-3}$ & $1.167$ \\
$0.025$      & $1.380\times10^{-3}$ & $0.961$
           & $9.047\times10^{-4}$ & $1.269$ \\
$0.0125$     & $6.992\times10^{-4}$ & $0.981$
           & $3.588\times10^{-4}$ & $1.334$ \\
$0.00625$    & $3.656\times10^{-4}$ & $0.936$
           & $1.447\times10^{-4}$ & $1.310$ \\
\hline\hline

\multicolumn{5}{c}{\textbf{$t = T = 872.8$ s}} \\
\hline
$0.1$        & $1.965\times10^{-2}$ & --
           & $1.750\times10^{-2}$ & -- \\
$0.05$       & $1.155\times10^{-2}$ & $0.767$
           & $8.910\times10^{-3}$ & $0.974$ \\
$0.025$      & $6.314\times10^{-3}$ & $0.871$
           & $3.859\times10^{-3}$ & $1.207$ \\
$0.0125$     & $3.315\times10^{-3}$ & $0.930$
           & $1.468\times10^{-3}$ & $1.395$ \\
$0.00625$    & $1.725\times10^{-3}$ & $0.942$
           & $5.157\times10^{-4}$ & $1.509$ \\
\hline
\end{tabular}
\caption{Two-way coupling (T3), $\Delta t\propto h$ at $\mathrm{Co}=0.5$: discrete $L^2$ error of the pressure head $\psi$ against the manufactured solution.}
\label{tab:mms_t3_psi}
\end{table}

\begin{table}[t]
\centering
\begin{tabular}{c|cc|cc}
\hline
\multirow{2}{*}{$h_e$ [m]}
& \multicolumn{2}{c|}{\textbf{Low Order}}
& \multicolumn{2}{c}{\textbf{FCT}} \\
\cline{2-5}
& $\|e\|_{L^2}$ & Order
& $\|e\|_{L^2}$ & Order \\
\hline
\hline

\multicolumn{5}{c}{\textbf{$t = T/4 = 218.2$ s}} \\
\hline
$0.1$        & $1.786\times10^{-2}$ & --
           & $1.084\times10^{-2}$ & -- \\
$0.05$       & $9.863\times10^{-3}$ & $0.857$
           & $4.675\times10^{-3}$ & $1.214$ \\
$0.025$      & $4.922\times10^{-3}$ & $1.003$
           & $2.081\times10^{-3}$ & $1.168$ \\
$0.0125$     & $2.119\times10^{-3}$ & $1.216$
           & $9.769\times10^{-4}$ & $1.091$ \\
$0.00625$    & $8.002\times10^{-4}$ & $1.405$
           & $4.748\times10^{-4}$ & $1.041$ \\
\hline\hline

\multicolumn{5}{c}{\textbf{$t = T/2 = 436.4$ s}} \\
\hline
$0.1$        & $2.342\times10^{-2}$ & --
           & $1.445\times10^{-2}$ & -- \\
$0.05$       & $1.296\times10^{-2}$ & $0.853$
           & $6.656\times10^{-3}$ & $1.119$ \\
$0.025$      & $6.475\times10^{-3}$ & $1.001$
           & $3.113\times10^{-3}$ & $1.097$ \\
$0.0125$     & $2.739\times10^{-3}$ & $1.241$
           & $1.516\times10^{-3}$ & $1.038$ \\
$0.00625$    & $1.115\times10^{-3}$ & $1.297$
           & $7.567\times10^{-4}$ & $1.002$ \\
\hline\hline

\multicolumn{5}{c}{\textbf{$t = T = 872.8$ s}} \\
\hline
$0.1$        & $1.856\times10^{-2}$ & --
           & $1.658\times10^{-2}$ & -- \\
$0.05$       & $9.922\times10^{-3}$ & $0.903$
           & $7.618\times10^{-3}$ & $1.122$ \\
$0.025$      & $4.951\times10^{-3}$ & $1.003$
           & $3.615\times10^{-3}$ & $1.075$ \\
$0.0125$     & $2.090\times10^{-3}$ & $1.244$
           & $1.804\times10^{-3}$ & $1.003$ \\
$0.00625$    & $8.741\times10^{-4}$ & $1.258$
           & $9.253\times10^{-4}$ & $0.963$ \\
\hline
\end{tabular}
\caption{Two-way coupling (T3), $\Delta t\propto h$ at $\mathrm{Co}=0.5$: discrete $L^2$ error of the concentration $c$ against the manufactured solution.}
\label{tab:mms_t3_c}
\end{table}

\begin{table}[t]
  \centering
  \caption{Coupling strength: $\|e\|_{L^2}(c)$ of T3 with FCT on $N=40$ at $t=T$, for
    two time steps, and their ratio.}
  \label{tab:mms_sweep}
  \small
  \begin{tabular}{rr ccc}
    \hline
    $\epsilon$ & $\omega_{\max}$ & $\Delta t=10.9$~s & $\Delta t/2$ & ratio \\
    \hline
    0     & 0    & $3.92\times10^{-3}$ & $2.00\times10^{-3}$ & 1.96 \\
    0.1   & 0    & $3.76\times10^{-3}$ & $1.96\times10^{-3}$ & 1.92 \\
    0.235 & 0    & $3.61\times10^{-3}$ & $1.97\times10^{-3}$ & 1.84 \\
    0.4   & 0    & $3.55\times10^{-3}$ & $2.09\times10^{-3}$ & 1.70 \\
    0.235 & 0.1  & $3.65\times10^{-3}$ & $1.94\times10^{-3}$ & 1.89 \\
    0.235 & 0.25 & $3.66\times10^{-3}$ & $1.94\times10^{-3}$ & 1.89 \\
    \hline
  \end{tabular}
\end{table}
